%==========================================================================================
\section{Diffraction-image scripts}\label{sec:images}
%==========================================================================================

Two command-line scripts process the detector frames of a BLISS mesh scan, typically the
Laue patterns of a microdiffraction map \citep{Ulrich2011}. Both read the scan through
\fn{bliss_hdf5_reader.BlissMeshScan} and parallelise the reading over processes.

%------------------------------------------------------------------------------------------
\subsection{\texttt{reduce\_image\_series}: minimum, maximum and mean images}
%------------------------------------------------------------------------------------------

\fn{python reduce_image_series.py master.h5 -s 1 -o OUTPUT_DIR [--image_suffix .h5] [-c 50]
[-w N]} computes, pixel by pixel, the minimum, maximum and mean over every frame of a scan,
and saves them in \fn{<master>_scanNNNN_reduced.h5}. The maximum image shows every Laue spot
that appears anywhere in the map; the minimum image is the background common to all frames
(useful to subtract, and to find hot or dead pixels); the mean is the average pattern.

The reduction is a map--reduce: the file list is split into chunks of \fn{chunk_size} files;
each worker returns the minimum, maximum and sum of its chunk and the number of valid frames;
the partial results are combined as they arrive (minimum and maximum are associative, sums
add), and the mean is the total sum over the number of frames. Sums are accumulated in
float64 to avoid integer overflow over thousands of frames.

\begin{issue}
The script iterates over image \emph{files} (\fn{scan.image_path_list}) and
\fn{_read_frame_array} returns the \emph{first} frame of each file. This is correct when the
detector saves one frame per file; when files hold several frames, only one frame per file is
reduced, and the \fn{processed_frames} attribute (the number of files) overstates the frames
used. The frame-level access of the reader (\fn{iter_frames}, as in
\fn{make_dynamic_mosaic}) would remove this restriction. The scan must also be
two-dimensional (\fn{BlissMeshScan}).
\end{issue}

\pycode{reduce_image_series}{_header}

\api[\texttt{parse\_args}]{\texttt{parse\_args()}}
Command-line arguments: the master file, \fn{-s/--scan} (scan number), \fn{-o/--output_dir},
\fn{--image_suffix} (default \fn{.h5}), \fn{-c/--chunk_size} (files per task, default 50),
\fn{-w/--workers} (default: all cores).
\pycode{reduce_image_series}{parse_args}

\api[\texttt{\_read\_frame\_array}]{\texttt{\_read\_frame\_array(filepath)}}
Reads the image of a detector file: first the usual Lima/BLISS dataset paths
(\fn{/entry_0000/measurement/data}, \fn{.../instrument/eiger4m/data},
\fn{.../CRGIF/eiger4m/data}), else the first dataset of two or more dimensions found under
\fn{/entry_0000}; for a 3D dataset the first frame is returned. Raises \fn{KeyError} if no
image is found.
\pycode{reduce_image_series}{_read_frame_array}

\api[\texttt{\_worker\_process\_files}]{\texttt{\_worker\_process\_files(file\_paths)}}
The map step: running minimum, maximum and float64 sum over a list of files, skipping
unreadable files with an error message. Returns \fn{(min, max, sum, n_valid)}, with
\fn{None} arrays if no file could be read.
\pycode{reduce_image_series}{_worker_process_files}

\api[\texttt{compute\_scan\_reductions}]{\texttt{compute\_scan\_reductions(image\_files, chunk\_size=50, num\_workers=None)}}
The reduce step, with a \fn{multiprocessing.Pool} and a progress bar. Returns
\fn{(min_image, max_image, mean_image)}; raises if the list is empty or no frame could be
read.
\pycode{reduce_image_series}{compute_scan_reductions}

\api[\texttt{save\_reduced\_hdf5}, \texttt{main}]{\texttt{save\_reduced\_hdf5(output\_path, global\_min, global\_max, global\_mean, metadata)}, \texttt{main()}}
Writes the three images (gzip level 4) and the metadata (master file, scan number and title,
number of files) as root attributes. \fn{main} wires the steps together.
\pycode{reduce_image_series}{save_reduced_hdf5}
\pycode{reduce_image_series}{main}
\pycode{reduce_image_series}{_tail}

%------------------------------------------------------------------------------------------
\subsection{\texttt{make\_dynamic\_mosaic}: region-of-interest maps}
%------------------------------------------------------------------------------------------

\fn{python make_dynamic_mosaic.py master.h5 -s SCAN -rx X -ry Y [-w HWS] [-o DIR]} follows one
region of interest (ROI) of the detector, a square of side $2\,\mathrm{hws}+1$ pixels centred
on $(X, Y)$---usually around one Laue spot---through every frame of a mesh scan. It produces:
\begin{itemize}[nosep]
\item the \textbf{mosaic}: the ROIs tiled in the order of the scan grid, a real-space image of
  how the spot changes across the sample;
\item the \textbf{entropy} map: the Shannon entropy \citep{Shannon1948} of the normalised ROI
  intensity, $H = -\sum_k p_k\ln p_k$ with $p_k = I_k/\sum I$. A sharp spot concentrates the
  intensity in a few pixels (low $H$); a diffuse or split spot, or background only, spreads it
  (high $H$, at most $\ln N_\mathrm{pixels}$ for a flat ROI);
\item the \textbf{max-to-mean} ratio of each ROI, another sharpness indicator;
\item the \textbf{centre of mass} of each ROI, $\bar x = \sum_k x_kI_k/\sum_k I_k$, as a
  displacement from its mean over the map: the motion of a Laue spot on the detector follows
  the local lattice orientation (and the deviatoric strain), so this vector field maps
  rotations across the sample;
\item the \textbf{total intensity} of the ROI, normalised by the monitor;
\item the sharpest and the most diffuse ROI.
\end{itemize}
The arrays are saved in \fn{computed_metrics_scanNNNN.h5} and a six-panel figure in physical
motor coordinates.

\begin{hyp}
\begin{itemize}[nosep,leftmargin=1em]
\item The ROI is fixed in detector pixels and must lie entirely inside the frame (a negative
  start index gives an empty slice); the spot must stay inside it over the whole map.
\item No background is subtracted: a flat background raises the entropy and pulls the centre
  of mass towards the ROI centre.
\item Displacements are in detector pixels, not converted into angles or strain.
\item Points never recorded (aborted scan) become empty blocks, and their motor positions
  are extrapolated from a plane fitted to the recorded ones.
\end{itemize}
\end{hyp}

\pycode{make_dynamic_mosaic}{_header}

\api[\texttt{parse\_args}]{\texttt{parse\_args()}}
The master file, \fn{-s/--scan}, \fn{-rx/--roi_x}, \fn{-ry/--roi_y}, \fn{-w/--hws} (half
window, default 5), \fn{--image_suffix}, \fn{-o/--output_dir}.
\pycode{make_dynamic_mosaic}{parse_args}

\api[\texttt{read\_image}, \texttt{roi\_slices}, \texttt{get\_image\_roi}]{\texttt{read\_image(filepath)}, \texttt{roi\_slices(roi\_x, roi\_y, hws)}, \texttt{get\_image\_roi(filepath, roi\_x, roi\_y, hws)}}
Fallback readers for scans whose master file gives no direct access to the frames (TIFF
files, or HDF5 files at \fn{/entry_0000/CRGIF/eiger4m/data}, first frame); \fn{roi_slices}
returns the (row, column) slices of the ROI.
\pycode{make_dynamic_mosaic}{read_image}
\pycode{make_dynamic_mosaic}{roi_slices}
\pycode{make_dynamic_mosaic}{get_image_roi}

\api[\texttt{\_init\_worker}, \texttt{\_read\_roi\_batch}, \texttt{read\_rois}]{\texttt{\_init\_worker(scan)}, \texttt{\_read\_roi\_batch(indices, roi)}, \texttt{read\_rois(scan, roi, batch=50)}}
Frames are read through the master file with \fn{BlissScan.iter_frames}, which reads only the
ROI from disk. The scan object is sent once to each worker process (pool initialiser, global
\fn{_WORKER_SCAN}), and each task reads a batch of 50 frames with the file opened once. Returns
the ROIs of every recorded frame in acquisition order.
\pycode{make_dynamic_mosaic}{_init_worker}
\pycode{make_dynamic_mosaic}{_read_roi_batch}
\pycode{make_dynamic_mosaic}{read_rois}

\api[\texttt{fill\_missing\_positions}, \texttt{to\_scan\_grid}]{\texttt{fill\_missing\_positions(grid)}, \texttt{to\_scan\_grid(rois, index\_grid)}}
\fn{fill_missing_positions} replaces the NaN positions of unrecorded points by a plane
$a + b\,i + c\,j$ fitted (least squares) to the recorded positions, since plotting functions
refuse NaN coordinates. \fn{to_scan_grid} orders the ROIs like the grid through the scan's
\fn{index_grid} (which handles back-and-forth meshes), with zero blocks for missing points.
\pycode{make_dynamic_mosaic}{fill_missing_positions}
\pycode{make_dynamic_mosaic}{to_scan_grid}

\api[\texttt{compute\_pipeline\_data}]{\texttt{compute\_pipeline\_data(array\_list, grid\_shape)}}
All metrics in vectorised form. The mosaic is built by reshaping the stack of blocks to
$(n_\mathrm{rows}, n_\mathrm{cols}, h, w)$, swapping the two middle axes and reshaping to
$(n_\mathrm{rows}h, n_\mathrm{cols}w)$. Empty blocks get zero entropy, max-to-mean and
displacement, and are ignored when looking for the sharpest block. Returns a dictionary of
arrays.
\pycode{make_dynamic_mosaic}{compute_pipeline_data}

\api[\texttt{save\_metrics\_to\_hdf5}, \texttt{plot\_and\_save\_figure}]{\texttt{save\_metrics\_to\_hdf5(data, output\_path)}, \texttt{plot\_and\_save\_figure(data, x\_grid, y\_grid, mon\_grid, motor\_names, output\_path)}}
The HDF5 archive (gzip), and the figure: mosaic in physical extent, sharpest and most diffuse
ROI, entropy map, monitor-normalised intensity, and the centre-of-mass vector field (one arrow
in two on dense grids).
\pycode{make_dynamic_mosaic}{save_metrics_to_hdf5}
\pycode{make_dynamic_mosaic}{plot_and_save_figure}

\api[\texttt{main}]{\texttt{main()}}
Reads the scan and the monitor, chooses as horizontal axis the motor whose name contains
``x'' (transposing the grids if it is the slow motor, so that the mosaic is drawn with $x$
horizontal), reads the ROIs (through the master file, or file by file), computes, saves and
plots.
\pycode{make_dynamic_mosaic}{main}
\pycode{make_dynamic_mosaic}{_tail}
